\section{Capitolo 2: Funzioni e Programmazione Funzionale}

\begin{esercizio}{Es. 5: Pipeline di Composizione}
\textbf{TRACCIA:} Scrivere una funzione `pipeline(funzioni)` che accetta un array di funzioni $[f_1, f_2, \dots, f_n]$ e restituisce una nuova funzione che applica le $f$ in sequenza all'input.
$result = f_n(\dots f_2(f_1(x))\dots)$

\textbf{Esempio:} `[x => x+1, x => x*2]` $\to$ Input 5 $\to$ (5+1)=6 $\to$ 6*2=12.
\end{esercizio}

\begin{jsbox}
function pipeline(funzioni) {
    return function(input) {
        let valore = input;
        for (let f of funzioni) {
            valore = f(valore); // Passo il risultato alla prossima
        }
        return valore;
    };
}
// Equivalente con reduce:
// const pipeline = fs => x => fs.reduce((acc, f) => f(acc), x);
\end{jsbox}

\begin{esercizio}{Es. 6: Filtro Asimmetrico (HOF)}
\textbf{TRACCIA:} Scrivere `partition(arr, predicato)` che restituisce un oggetto `{ pass: [...], fail: [...] }`.
In `pass` vanno gli elementi per cui il predicato è true, in `fail` gli altri.

\textbf{Vincolo:} Scorrere l'array una sola volta.
\end{esercizio}

\begin{jsbox}
function partition(arr, predicato) {
    let result = { pass: [], fail: [] };
    
    for (let el of arr) {
        if (predicato(el)) {
            result.pass.push(el);
        } else {
            result.fail.push(el);
        }
    }
    return result;
}
\end{jsbox}

\begin{esercizio}{Es. 7: Generatore di Contatori (Closure)}
\textbf{TRACCIA:} Scrivere `creaContatori(n)` che restituisce un array di `n` funzioni.
La funzione all'indice `i` deve ritornare `i` quando chiamata.
\textit{Attenzione all'errore classico del `var` nei cicli!}
\end{esercizio}

\begin{jsbox}
function creaContatori(n) {
    let funzioni = [];
    
    for (let i = 0; i < n; i++) {
        // Usando 'let' nel ciclo, creiamo uno scope per ogni iterazione.
        // La closure cattura la 'i' corretta.
        funzioni.push(function() {
            return i;
        });
    }
    return funzioni;
}

// Se usassimo 'var i', tutte ritornerebbero n!
\end{jsbox}

\begin{esercizio}{Es. 8: Map su Albero (Tree Map)}
\textbf{TRACCIA:} Scrivere una funzione ricorsiva `alberoMap(nodo, f)` che restituisce un \textbf{nuovo} albero strutturalmente identico, ma dove ogni valore `val` è trasformato dalla funzione `f`.
\end{esercizio}

\begin{jsbox}
function alberoMap(nodo, f) {
    if (!nodo) return null;
    
    return {
        val: f(nodo.val),           // Trasformo il valore
        sx: alberoMap(nodo.sx, f),  // Ricorro a sinistra
        dx: alberoMap(nodo.dx, f)   // Ricorro a destra
    };
}
\end{jsbox}

\begin{esercizio}{Es. 20: Memoization (Cache Semplice)}
\textbf{TRACCIA:} Scrivere una funzione `memoize(f)` che accetta una funzione (che prende 1 argomento stringa/numero) e ne restituisce una versione "con memoria". Se richiamata con lo stesso input, non ricalcola ma restituisce il valore salvato.
\end{esercizio}

\begin{jsbox}
function memoize(f) {
    let cache = {}; // La chiusura protegge la cache
    
    return function(n) {
        if (n in cache) {
            return cache[n]; // HIT
        } else {
            let result = f(n); // MISS - Calcolo
            cache[n] = result;
            return result;
        }
    };
}
\end{jsbox}

\begin{esercizio}{Es. 21: Currying (Somma a Catena)}
\textbf{TRACCIA:} Scrivere una funzione `somma(a)` che può essere chiamata così: `somma(2)(3)(4)` e restituisce la somma finale (es. supportando 3 chiamate fisse o tornando un oggetto con `valueOf`).
\textit{Versione Semplice (3 argomenti fissi):}
\end{esercizio}

\begin{jsbox}
function somma(a) {
    return function(b) {
        return function(c) {
            return a + b + c;
        };
    };
}
// somma(1)(2)(3) -> 6
\end{jsbox}

\begin{esercizio}{Es. 22: Implementazione Manuale di 'Every'}
\textbf{TRACCIA:} Scrivere `myEvery(arr, predicato)` che torna true se \textbf{TUTTI} gli elementi soddisfano il predicato.
\textit{Appena ne trovi uno falso, devi fermarti (short-circuit).}
\end{esercizio}

\begin{jsbox}
function myEvery(arr, predicato) {
    for (let el of arr) {
        if (!predicato(el)) {
            return false; // Trovato un contromodello
        }
    }
    return true; // Nessuno ha fallito
}
\end{jsbox}

\begin{esercizio}{Es. 47: Map tramite Reduce}
\textbf{TRACCIA:} Implementare `myMap(arr, f)` usando \textbf{solo} `reduce`.
\end{esercizio}

\begin{jsbox}
function myMap(arr, f) {
    return arr.reduce((acc, el) => {
        acc.push(f(el));
        return acc; // Ricorda di ritornare l'accumulatore!
    }, []); // Valore iniziale array vuoto
}
\end{jsbox}

\begin{esercizio}{Es. 48: Filter tramite Reduce}
\textbf{TRACCIA:} Implementare `myFilter(arr, pred)` usando \textbf{solo} `reduce`.
\end{esercizio}

\begin{jsbox}
function myFilter(arr, predicato) {
    return arr.reduce((acc, el) => {
        if (predicato(el)) {
            acc.push(el);
        }
        return acc;
    }, []);
}
\end{jsbox}
